\documentclass[10pt]{article}
\usepackage[T1]{fontenc}
\usepackage[utf8]{inputenc}
\usepackage[french]{babel}
\usepackage{newpxtext,newpxmath}
\usepackage[a4paper,top=1.5cm,bottom=1.35cm,left=1.65cm,right=1.65cm]{geometry}
\usepackage{microtype}
\usepackage{xcolor}
\usepackage{titlesec}
\usepackage[most]{tcolorbox}
\usepackage{enumitem}
\usepackage{array}
\usepackage{booktabs}
\usepackage{fancyhdr}

\definecolor{navy}{HTML}{1F3A5F}
\definecolor{grey}{HTML}{54606E}
\definecolor{green}{HTML}{2C6E49}
\definecolor{amber}{HTML}{9C6A16}
\definecolor{red}{HTML}{A6322E}
\definecolor{plum}{HTML}{7A4EA8}
\definecolor{boxgray}{HTML}{F3F5F8}
\definecolor{boxred}{HTML}{FCF2F1}
\definecolor{boxgreen}{HTML}{F1F7F3}
\definecolor{boxedge}{HTML}{D9E0EA}

\setlength{\parindent}{0pt}
\setlength{\parskip}{2.6pt}
\linespread{1.0}

\pagestyle{fancy}\fancyhf{}
\lhead{\small\color{grey}Scella}
\rhead{\small\color{grey}État des lieux -- 10 août 2026}
\cfoot{\small\color{grey}\thepage}
\renewcommand{\headrulewidth}{0.3pt}

\titleformat{\section}{\color{navy}\large\bfseries}{}{0pt}{}
\titlespacing{\section}{0pt}{9pt}{3pt}
\titleformat{\subsection}{\color{navy}\bfseries}{}{0pt}{}
\titlespacing{\subsection}{0pt}{7pt}{1.5pt}

\newtcolorbox{these}{colback=boxgray,colframe=boxedge,boxrule=0.5pt,arc=2.5pt,
  left=8pt,right=8pt,top=4.5pt,bottom=4.5pt}
\newtcolorbox{alerte}{colback=boxred,colframe=red!55,boxrule=0.8pt,arc=2.5pt,
  left=8pt,right=8pt,top=4.5pt,bottom=4.5pt}
\newtcolorbox{acquis}{colback=boxgreen,colframe=green!45,boxrule=0.8pt,arc=2.5pt,
  left=8pt,right=8pt,top=4.5pt,bottom=4.5pt}

\newcommand{\ok}{{\color{green}\textbf{FAIT}}}
\newcommand{\partiel}{{\color{amber}\textbf{PARTIEL}}}
\newcommand{\design}{{\color{grey}\textbf{CONÇU}}}
\newcommand{\nope}{{\color{red}\textbf{NON FAIT}}}
\newcommand{\sz}[1]{{\small\color{grey}\textsf{#1}}}
\newcommand{\code}[1]{\texttt{\small #1}}

\begin{document}

\begin{center}
{\Huge\bfseries\color{navy} Scella}\\[4pt]
{\large\color{grey}\itshape Où en est-on : état des lieux mesuré au code}\\[4pt]
{\small\color{grey} 10 août 2026 \; · \; branche \code{couche2-m0-frontier}, HEAD \code{0077fcb} \; · \; 66 crates}
\end{center}
\vspace{2pt}

\begin{these}
\textbf{Le résultat central.} Le méta-essaim est passé, en une journée, de \emph{coordonner} à
\textbf{attester} : une ronde servie rend désormais un reçu qu'un tiers re-prouve hors ligne, qui
\textbf{nomme} ses agents et sa règle, et dont une \textbf{troncature} se voit \textbf{sans la clé}.
C'est le premier assemblage complet et montrable du dépôt. Mais la mesure fait apparaître un fait que
personne n'avait écrit : \emph{toutes} ces capacités sont \textbf{opt-in}, et \textbf{aucun profil de
déploiement ne les allume ensemble}. Un build par défaut de Scella ne sert donc \textbf{aucun reçu} --
et la seule configuration cohérente qui existe aujourd'hui est un script de démonstration.
\end{these}

\section{1 · Méthode : ce que « mesuré » veut dire ici}

Aucune affirmation de ce document ne vient d'un document de conception ni d'un souvenir de session.
Chaque chiffre, chaque « fait » et chaque « non fait » a été relu dans le code de ce matin, à
l'emplacement indiqué. C'est la discipline habituelle du projet, et elle a de nouveau payé : trois
énoncés que je tenais pour acquis se sont révélés inexacts au contrôle, et ils sont corrigés
\emph{à leur place} plutôt qu'en silence.

Le premier concerne l'ancrage tiers. J'avais avancé que RFC 3161 « n'apparaît qu'en commentaire ».
C'est faux : le crate \code{forensic-export} porte déjà un \code{AnchorMethod::Rfc3161}, une structure
\code{AnchorProof} et un \code{verify\_anchor} qui exécute le contrôle décisif -- la liaison de racine
(\code{bound\_root == manifest\_root}). Ce qui est délégué, et le module le dit lui-même, c'est le
\emph{parsing cryptographique du jeton}. La conséquence est favorable et sera reprise au §5 : le jalon
3b n'est pas un travail à inventer, c'est un raccordement.

Le deuxième concerne l'état de la branche. « Jamais fusionnée » était inexact : \code{main} pointe sur
\code{d2613fa}, un merge de \code{couche2-m0-frontier} daté du \textbf{12 juillet}. La branche a donc
bien été intégrée une fois ; ce sont les \textbf{31 commits accumulés depuis} qui ne l'ont pas été.

Le troisième concerne le périmètre : ce document couvre la Couche 2, la Couche 3 et les deux piles de
reçu. Il ne couvre \emph{pas} le Fil B ni la couche nano, qui n'ont pas été relus dans cet arc. Le §6
le dit explicitement plutôt que de laisser le silence passer pour une validation.

\section{2 · L'arc de la journée : de « coordonner » à « attester »}

Huit commits ont été poussés aujourd'hui, mais ils ne forment qu'un seul mouvement, et c'est ainsi
qu'il faut les lire. Avant ce matin, le méta-essaim faisait converger plusieurs participants et
n'en laissait \textbf{aucune trace opposable} : la plomberie était auditable, la ronde ne produisait
rien. La question qu'un évaluateur pose en premier -- « montrez-moi ce que la machine a
\emph{produit}, pas comment elle est faite » -- restait sans réponse.

Elle en a une maintenant, et cette réponse satisfait trois propriétés qu'il faut distinguer, parce
qu'elles se démontrent séparément et qu'aucune n'implique les autres.

\renewcommand{\arraystretch}{1.2}
\begin{center}\small
\begin{tabular}{@{}l >{\raggedright\arraybackslash}p{7.4cm} >{\raggedright\arraybackslash}p{6.1cm}@{}}
\toprule
\textbf{Jalon} & \textbf{Propriété obtenue} & \textbf{Ce qui la démontre} \\
\midrule
\textbf{M2} & Un tiers re-prouve le reçu \emph{hors ligne}, sans faire confiance au serveur &
  \code{cargo tree -e normal} montre le vérifieur \textbf{seul} : zéro dépendance \\
\textbf{v4} & Le reçu \textbf{nomme} ses acteurs, leurs versions déclarées et sa règle d'agrégation &
  Renommer un agent, permuter deux colonnes ou changer $\alpha$ : la re-preuve \textbf{casse} \\
\textbf{3a} & Une \textbf{troncature} est détectable \textbf{sans aucune clé} &
  Le bundle tronqué passe la re-preuve avec la clé, et l'ancre le \textbf{refuse} sans \\
\bottomrule
\end{tabular}
\end{center}

\subsection{2.1 · Pourquoi l'ancre était indispensable, et non décorative}

Le point mérite d'être posé lentement, parce qu'il est contre-intuitif et qu'il est au cœur de ce que
le projet appelle un \emph{vert malhonnête} -- un système qui affiche OK en faisant autre chose.

Une chaîne HMAC prouve que ce qui est là n'a pas bougé. Elle ne prouve pas que \emph{tout} est encore
là. La raison est structurelle : \textbf{le préfixe d'une chaîne HMAC valide est lui-même une chaîne
HMAC valide}. Retirer les derniers pas d'une ronde et « corriger » le manifeste produit donc un bundle
qui se re-prouve \textbf{parfaitement} -- rien, à l'intérieur, ne trahit la coupure. Un opérateur
mécontent d'une ronde peut ainsi en publier la moitié qui l'arrange, et le vérifieur dira vert.

Seule une \textbf{racine publiée à l'avance}, accompagnée du nombre de pas, expose la coupure. Et le
contrôle qui l'expose ne demande \textbf{aucune clé} : c'est une comparaison de deux chaînes de
caractères contre un journal append-only. Le jalon 3a a donc ajouté une propriété que la cryptographie
seule ne pouvait pas donner, et la démo l'exhibe par un « vert gênant » -- elle montre d'abord le
bundle tronqué \emph{passer} la re-preuve, avant de le faire tomber sur l'ancre. Un contrôle qui ne
peut pas échouer ne prouve rien ; le montrer échouer est ce qui rend la démonstration crédible.

\subsection{2.2 · Ce que la démonstration montre, et ce qu'elle ne montre pas}

L'ensemble tient en une commande, en deux modes. Sans drapeau, trois participants \emph{mock} : le
mécanisme est démontré, mais la colonne de configuration du reçu est vide -- un mock n'a aucune
version à attester, et c'est la première chose qu'un lecteur attentif remarque. Avec \code{--real},
trois participants \textbf{efe-core} réels, chacun déclaré par un bloc de contrôleur dont le SHA-256
est scellé, chacun passant au démarrage le canari H1 de non-expansivité en \emph{fail-closed}. Le reçu
affiche alors trois \code{config\_hash} distincts et réels, ce qui autorise enfin la phrase qu'un
acheteur conformité veut entendre : \emph{cet agent, dans cette version précise, a contribué ceci}.

Deux limites doivent accompagner toute présentation. Les participants vivent dans le processus de la
démo : le transport gRPC est réel, l'isolation système d'exploitation ne l'est pas -- choix de
fiabilité assumé, l'orchestration multi-processus du dépôt étant instable sous charge, et une démo qui
échoue une fois sur cinq ne démontre rien. Et surtout, le reçu atteste l'\textbf{intégrité} de la
suite de pas, \textbf{jamais la correction} de la décision de consensus. Il est \emph{tamper-evident
pour un détenteur de clé}, pas \emph{trustless pour un tiers} : un détenteur peut fabriquer une chaîne
cohérente. C'est le cadrage constant du projet, et il est répété jusque dans le corps de la réponse
HTTP, à l'endroit précis où un lecteur pressé conclurait « vérifié = correct ».

\section{3 · Le fait structurel : tout est opt-in, aucun profil ne l'allume}

C'est le résultat le plus important de cet état des lieux, et il n'était visible qu'en regardant les
\emph{défauts} plutôt que les capacités. Chaque crate servi part désactivé.

\renewcommand{\arraystretch}{1.1}
\begin{center}\small
\begin{tabular}{@{}l l >{\raggedright\arraybackslash}p{9.3cm}@{}}
\toprule
\textbf{Crate servi} & \textbf{\code{default}} & \textbf{Ce que le défaut laisse dehors} \\
\midrule
\code{efe-swarm-gateway}     & \code{[]}      & la frontière LLM, le garde d'egress, le journal d'ancres \\
\code{efe-meta-swarm-server} & \code{["tls"]} & le moteur efe-core réel (\code{core-integration}) \\
\code{efe-server}            & \code{[]}      & la façade \code{/v1/consensus} et son client méta-essaim \\
\code{efe-meta-swarm}        & \code{[]}      & les chemins optionnels du moteur \\
\bottomrule
\end{tabular}
\end{center}

À quoi s'ajoutent sept interrupteurs d'environnement, répartis sur trois crates, qu'il faut connaître
et poser \emph{ensemble} pour obtenir le comportement que décrivent les documents de présentation.

\begin{center}\small
\begin{tabular}{@{}l >{\raggedright\arraybackslash}p{9.9cm}@{}}
\toprule
\textbf{Interrupteur} & \textbf{Ce qu'il commande} \\
\midrule
\code{EFE\_META\_SWARM\_ENABLED}            & l'existence même du sous-système méta-essaim \\
\code{EFE\_META\_PROVENANCE\_TENANT}        & le scellement du reçu ; absent $\Rightarrow$ \textbf{aucun reçu} \\
\code{EFE\_NANO\_AUDIT\_HMAC\_MASTER}       & le master dont la clé tenant est dérivée \\
\code{EFE\_META\_PROVENANCE\_ALLOW\_DEV\_KEY} & l'autorisation explicite d'une clé de développement \\
\code{EFE\_SWARM\_ANCHOR\_LOG}              & le journal d'ancres -- sans lui, pas de détection de troncature \\
\code{EFE\_LLM\_FRONTIER\_ENABLED}          & les appels aux modèles frontière (Couche 2) \\
\code{EFE\_LLM\_ACCEPT\_CLEARTEXT\_EGRESS}  & la reconnaissance explicite d'une sortie en clair \\
\bottomrule
\end{tabular}
\end{center}

\begin{alerte}
\textbf{Le constat.} La recherche a été faite et elle est revenue vide : \textbf{aucun} des fichiers de
déploiement du dépôt (\code{.env.slm}, \code{.env.template}, \code{.envrc}, \code{poc.env},
\code{docker/Dockerfile.prod}) ne pose un seul des quatre interrupteurs du reçu, et \textbf{aucun
document hors code} ne mentionne \code{EFE\_META\_PROVENANCE\_TENANT}. Conséquence exacte : un build par
défaut de Scella ne sert \textbf{aucun reçu} -- pas un reçu faible, aucun. La racine part vide, ce qui
est le signal honnête câblé à dessein, mais le fait demeure : « Scella fait X » n'est vrai que sous une
configuration qui n'existe nulle part comme artefact nommé, sauf dans \code{scripts/demo-consensus-receipt.sh}.
\end{alerte}

Il faut être juste sur l'origine de cette situation : elle ne vient pas d'un oubli mais d'une
discipline. Le \emph{fail-closed} est une vertu, et l'alternative -- servir un reçu sous clé de
développement sans que personne l'ait demandé -- serait exactement la malhonnêteté que tout le
dispositif existe pour empêcher. Le manque n'est donc pas dans le code, il est dans ce qui l'accompagne :
il manque \textbf{un profil nommé}, disant pour chaque interrupteur ce qu'il achète et ce que son
absence coûte. C'est une page, et c'est aujourd'hui l'écart le plus rentable du dépôt.

\section{4 · Le delta contre le rapport du 5 août}

Le dépôt dispose d'un point de comparaison daté et écrit : l'analyse composant par composant du
5 août, dont le verdict d'ouverture était \emph{« presque tout est construit ; presque rien n'est
servi »}. Cinq jours plus tard, quatre de ses lignes ont changé -- et sa propre alerte centrale est
fermée.

\renewcommand{\arraystretch}{1.15}
\begin{center}\small
\begin{tabular}{@{}>{\raggedright\arraybackslash}p{3.5cm} l >{\raggedright\arraybackslash}p{8.6cm}@{}}
\toprule
\textbf{Ligne du 5 août} & \textbf{Aujourd'hui} & \textbf{Ce qui a bougé, et où le vérifier} \\
\midrule
Le méta-essaim \newline \sz{NON SERVI · XL} & \partiel &
  M0, M1 et M2 faits ; moteur réel vérifié (canari H1, e2e DoD\#1). Restent M3 et M4 \\
Le reçu, pile A \newline \sz{non résistant à la troncature} & \ok &
  \code{routes/verify.rs:74--78} : \code{root\_matches} existe et entre dans le verdict \\
Le reçu, pile B \newline \sz{aucun producteur} & \ok &
  \code{GET /v1/swarm/bundle/:id} monté dans \code{efe-swarm-gateway/src/lib.rs:133} \\
Couche 2 servie \newline \sz{PII fail-open, budget no-op} & \ok &
  commits \code{a85f178} puis \code{af1e105} -- dont un identifiant qui ne corrélait rien \\
\bottomrule
\end{tabular}
\end{center}

La ligne « pile A » mérite un mot, parce que c'était l'alerte encadrée en rouge du 5 août : la route
qui portait le verdict public recalculait la racine et \textbf{ne la comparait jamais} à la racine
scellée, alors que la comparaison correcte existait un fichier plus loin. Un tiers disposant de
l'écriture en base et d'\emph{aucune clé} obtenait \code{verified: true} sur un run tronqué. La
comparaison est là aujourd'hui. Subsiste un \textbf{résidu assumé et documenté} : un run dont la
racine n'a pas été persistée rend \code{root\_matches: None} et conserve son verdict -- choix de
non-accusation rétroactive, défendable, mais qui doit être dit plutôt que découvert.

\section{5 · Les maillons faibles d'aujourd'hui}

Ils ont bougé. Le reçu est solide ; ses limites sont nommées, et toutes ont été vérifiées comme
\emph{non faites} plutôt que supposées telles.

\subsection{5.1 · M4 -- le \code{config\_hash} est \emph{déclaré}, pas \emph{rapporté}}

C'est la limite la plus subtile et la plus importante à énoncer correctement. Le reçu v4 scelle la
configuration que le \textbf{manifeste demandait}, pas celle que le participant a \textbf{dit servir}.
Un participant qui répondrait en servant autre chose que sa ligne déclarée ne serait pas démasqué par
le reçu. La pièce manquante est le certificat porté sur le fil : \code{client/remote\_participant.rs:404}
retourne encore \code{Ok(NonExpansiveCertificate::Stub)}, un marqueur de phase 1 qui, comme son
module le dit, n'effectue \textbf{aucun contrôle}. Tant qu'il est là, le mode strict ne peut pas être
activé sur l'\code{Iterate} servi. \sz{Taille M.}

\subsection{5.2 · 3b -- l'ancre est détenue par l'opérateur}

Le journal posé aujourd'hui supprime la réécriture \emph{rétroactive} : on ne peut plus republier une
version tronquée d'une ronde déjà ancrée. Il ne supprime pas la confiance dans l'opérateur, qui reste
maître de son journal. La garantie tierce -- horodatage RFC 3161, témoin externe, log de transparence --
est le jalon suivant.

C'est ici que la correction annoncée au §1 change l'estimation. Le crate \code{forensic-export} porte
déjà l'\code{AnchorProof} structurée, l'énumération des méthodes (\code{rfc3161},
\code{transparency-log}, \code{regulator-deposit}, \code{declared}) et le \code{verify\_anchor} qui
exécute le contrôle décisif : la preuve commet-elle bien \textbf{ce} DAG. Ce qui manque du côté
Couche 3 n'est donc pas une invention mais un \textbf{raccordement} : émettre une \code{AnchorProof}
plutôt qu'une ligne de journal, et brancher la vérification du jeton chez l'émetteur. \sz{Taille M,
revue à la baisse.}

\subsection{5.3 · M3 -- aucune résistance aux valeurs aberrantes}

Le contrôle est net : zéro opérateur de médiane, de moyenne élaguée ou de type Krum dans le moteur.
L'agrégation servie est uniforme, donc un participant défaillant ou hostile déplace le consensus à
proportion de son poids. Ce jalon dépendait de v4 et ne pouvait pas être fait avant : une médiane
élaguée sur des votes \emph{positionnels} n'est pas auditable -- pour justifier d'avoir écarté une
contribution, il faut pouvoir la \textbf{nommer}. Le déblocage est acquis depuis ce matin. \sz{Taille L.}

\subsection{5.4 · Un arbitrage laissé ouvert}

Le DoD\#2 demande qu'une topologie dynamique fasse \textbf{échouer} la requête ; le code journalise un
avertissement et sert. Ce n'est pas un bug -- c'est une décision non tranchée, et elle doit l'être
explicitement, dans un sens ou dans l'autre, plutôt que rester un écart entre un document et un
comportement. \sz{Taille XS.}

\section{6 · Portée : ce que cet état des lieux ne couvre pas}

Cette analyse a relu la Couche 2 (conformité), la Couche 3 (conformité \emph{et} actualité des
documents) et les deux piles de reçu. Le Fil P l'avait été précédemment. N'ont \textbf{pas} été
audités dans cet arc : le \textbf{Fil B} (le relais SLM entre modèles frontière) et la \textbf{couche
nano}. Leur silence ici n'est pas un satisfecit.

Une sonde a par ailleurs été passée sur les documents de la Couche 2, avec la signature qui avait
révélé huit péremptions côté Couche 3 (« n'existe pas », « 0 hit », « non câblé »). Trois occurrences,
\textbf{toutes disculpées} après lecture : un récit de piège corrigé, une table de corrections, un
extrait de code. Ce n'est pas un balayage complet -- c'est une sonde, et elle est revenue propre.

\section{7 · Position, et les deux choses qui manquent}

La branche \code{couche2-m0-frontier} porte \textbf{31 commits} non intégrés depuis le merge du
12 juillet. L'arbre de travail est propre, aux fichiers d'une session concurrente près (quatre sources
et deux PDF, intacts). Le dépôt compte 66 crates ; la chaîne du reçu -- moteur servi, bibliothèque de
vérification, binaire tiers -- est couverte par plusieurs centaines de tests dont le détail est mesuré
séparément.

\begin{acquis}
\textbf{Pour l'objectif réel, la capacité n'est plus le goulot.} L'artefact de conversion existe, il
tourne en une commande, et il survit à la première question dure -- \emph{quel agent, dans quelle
version, combiné selon quelle règle} -- à laquelle un reçu positionnel ne pouvait pas répondre. C'était
précisément l'objet des deux derniers jalons.
\end{acquis}

Ce qui manque tient en deux choses, et aucune des deux n'est technique.

\textbf{Le profil.} Sept interrupteurs et deux features qui font que Scella fait ce que ses documents
de présentation décrivent, écrits une fois, avec le coût de chaque absence. Sans lui, la seule
configuration cohérente qui existe est un script de démonstration -- et c'est une phrase qu'on ne veut
pas avoir à prononcer devant un évaluateur.

\textbf{La branche.} Trente-et-un commits qui ne sont visibles que si l'on sait quelle branche
regarder. Quiconque clone \code{main} aujourd'hui ne voit ni M2, ni v4, ni l'ancrage, ni la démo, ni le
job d'intégration continue qui défend le DoD\#1.

Le reste -- M4, 3b, M3 -- est du travail de qualité qui renforcera des propriétés que personne ne verra
dans une réunion de trente minutes.

\section{Annexe · Comment re-vérifier ce rapport}

Un état des lieux qui demande qu'on le croie sur parole a le même défaut que le \emph{vert malhonnête}
qu'il dénonce. Chaque affirmation du document est donc reproductible en une commande ou à une ligne de
code précise, depuis la racine du dépôt. Les affirmations \textbf{négatives} -- « M3 n'existe pas » --
sont les plus importantes à pouvoir contrôler, parce que ce sont celles qu'un lecteur ne peut pas
vérifier en lisant : une recherche qui revient vide est une preuve, à condition d'être rejouable.

\renewcommand{\arraystretch}{1.22}
\begin{center}\small
\begin{tabular}{@{}>{\raggedright\arraybackslash}p{5.05cm} >{\raggedright\arraybackslash}p{11.1cm}@{}}
\toprule
\textbf{Affirmation} & \textbf{Ce qui la rejoue} \\
\midrule
66 crates & \code{ls crates | wc -l} \\
31 commits non intégrés & \code{git rev-list -{}-count main..HEAD} \\
\code{main} fusionnée le 12 juillet & \code{git log -1 -{}-format=\%ci main} \\
Les défauts sont vides &
  \code{grep -A1 '\textasciicircum\textbackslash\lbrack features\textbackslash\rbrack' \newline
  crates/efe-\{server,swarm-gateway,meta-swarm\{,-server\}\}/Cargo.toml} \\
Aucun fichier de déploiement ne pose les interrupteurs du reçu &
  \code{grep -nE 'META\_PROVENANCE|SWARM\_ANCHOR' .env.slm .env.template .envrc poc.env docker/Dockerfile.prod} \\
Aucun document hors code ne les mentionne &
  \code{grep -rl EFE\_META\_PROVENANCE . -{}-include='*.md' -{}-include='*.yml'} \\
M4 non fait & \code{client/remote\_participant.rs:404} $\rightarrow$ \code{Ok(NonExpansiveCertificate::Stub)} \\
M3 non fait & \code{grep -rniE 'trimmed\_mean|krum|coordinate\_wise\_median' crates/efe-meta-swarm/src} \\
3b : la structure existe déjà & \code{forensic-export/src/anchor.rs:162} $\rightarrow$ \code{pub fn verify\_anchor} \\
Pile A : la racine est comparée & \code{efe-swarm-gateway/src/routes/verify.rs:74--78} \\
Pile B : le producteur existe & \code{efe-swarm-gateway/src/lib.rs:133} \\
Le vérifieur est sans dépendance & \code{cargo tree -p verify-meta-mini -e normal -{}-prefix none} \\
L'ensemble se démontre en une commande & \code{bash scripts/demo-consensus-receipt.sh -{}-real} \\
\bottomrule
\end{tabular}
\end{center}

Deux avertissements sur cette annexe, pour qu'elle ne soit pas lue plus largement qu'elle ne vaut. Les
lignes qui portent sur des \emph{absences} valent à la date du document et pour le périmètre indiqué au
§6 : elles ne disent rien du Fil B ni de la couche nano. Et la dernière ligne exige la feature
\code{core-integration}, donc quelques minutes de compilation au premier passage -- ce qui est
précisément le genre de détail qu'il vaut mieux découvrir avant une démonstration qu'au milieu.

\vspace{4pt}
\begin{these}
\small\textbf{Rappel de cadrage, à conserver dans toute présentation.} Le reçu atteste l'\textbf{intégrité}
d'une suite de pas, jamais la \textbf{correction} d'une décision. Il est \emph{tamper-evident pour un
détenteur de clé}, pas \emph{trustless pour un tiers}. Et les participants du mode réel exécutent une
fusion latente Mode-1 numérique : ce ne sont pas des « vrais SLM ».
\end{these}

\end{document}
